\chapter{Ma Trận}
\hypertarget{localtoc\thechapter}{}
\localtoc

\begin{problem}
    Hàm nhập xuất ma trận các số nguyên.
\end{problem}
\begin{problem}
    Hàm nhập xuất ma trận các số thực.
\end{problem}
\begin{problem}
    Viết hàm tìm giá trị lớn nhất trong ma trận các số thực.
\end{problem}
\begin{problem}
    Viết hàm kiểm tra trong ma trận các số nguyên có tồn tại giá trị chẵn nhỏ hơn 2004 hay không?
\end{problem}
\begin{problem}
     Viết hàm tính trung bình cộng các số nguyên tố trong ma trận số nguyên
\end{problem}
\begin{problem}
    Viết hàm tính tổng các giá trị âm trong ma trận các số thực.
\end{problem}
\begin{problem}
    Viết hàm sắp xếp ma trận các số thực tăng dần từ trên xuống dưới và từ trái sang phải.
\end{problem}

\section{Kĩ thuật tính toán}
\begin{problem}
    Tính tổng các số dương trong ma trận các số thực
\end{problem}
\begin{problem}
    Tính tích các giá trị lẻ trong ma trận các số nguyên
\end{problem}
\begin{problem}
     Tính tổng các giá trị trên 1 dòng trong ma trận các số thực
\end{problem}
\begin{problem}
    Tính tích các giá trị dương trên 1 cột trong ma trận các số thực
\end{problem}
\begin{problem}
    Tính tổng các giá trị dương trên 1 dòng trong ma trận các số thực
\end{problem}
\begin{problem}
    Tính tích các số chẵn trên 1 cột trong ma trận các số nguyên
\end{problem}
\begin{problem}
    Tính trung bình cộng các số dương trong ma trận các số thực
\end{problem}
\begin{problem}
    Tính tổng các giá trị nằm trên biên của ma trận
\end{problem}
\begin{problem}
    Tính trung bình nhân các số dương trong ma trận các số thực
\end{problem}
\begin{problem}
    Hãy biến đổi ma trận bằng cách thay các giá trị âm bằng trị tuyệt đối của nó
\end{problem}
\begin{problem}
    Hãy biên đổi ma trận số thực bằng cách thay các giá trị bằng giá trị nguyên gần nó nhất
\end{problem}
\begin{problem}
    Tính tổng các giá trị trên 1 dòng của ma trận các số thực
\end{problem}
\begin{problem}
    Tính tổng các giá trị lẻ trên 1 cột của ma trận các số nguyên
\end{problem}
\begin{problem}
    Tính tổng các số hoàn thiện trong ma trận các số nguyên
\end{problem}



\section{Kỹ thuật đếm}
\begin{problem}
    Viết hàm đếm số lượng số dương trong ma trận các số thực
\end{problem}
\begin{problem}
    Đếm số lượng số nguyên tố trong ma trận các số nguyên
\end{problem}
\begin{problem}
    Đếm tần suất xuất hiện của 1 giá trị x trong ma trận các số thực
\end{problem}
\begin{problem}
    Đếm số chữ số trong ma trận các số nguyên dương
\end{problem}
\begin{problem}
    Đếm số lượng số dương trên 1 hàng trong ma trận các số thực
\end{problem}
\begin{problem}
    Đếm số lượng số hoàn thiện trên 1 hàng trong ma trận các số nguyên
\end{problem}
\begin{problem}
    Đếm số lượng số âm trên 1 cột trong ma trận các số thực
\end{problem}
\begin{problem}
    Đếm số lượng số dương trên biên trong ma trận các số thực
\end{problem}
\begin{problem}
    * Đếm số lượng phần tử cực đại trong ma trận các số thực. Một phần tử được gọi là cực đại khi nó lớn hơn các phần tử xung quanh
\end{problem}
\begin{problem}
    * Đếm số lượng phần tử cực trị trong ma trận các số thực. Một phần tử được gọi là cực trị khi nó lớn hớn các phần tử xung quanh hoặc nhỏ hơn các phần tử xung quanh
\end{problem}
\begin{problem}
    * Đếm số lượng giá trị phân biệt có trong ma trận các số thực.

    Chú ý: Nếu có k phần tử (k >= 1) trong ma trận bằng nhau thì ta chỉ tính là 1
\end{problem}
\begin{problem}
    (*): Tính tổng các phần tử cực trị trong ma trận các số thực. Một phần tử được gọi là cực trị khi nó lớn hớn các phần tử xung quanh hoặc nhỏ hơn các phần tử xung quanh
\end{problem}
\begin{problem}
    (*): Đếm số lượng giá trị “Hoàng Hậu” trên ma trận. Một phần tử được gọi là Hoàng Hậu khi nó lớn nhất trên dòng, trên cột và 2 đường chéo đi qua nó
\end{problem}
\begin{problem}
    (*): Đếm số lượng giá trị “Yên Ngựa” trên ma trận. Một phần tử được gọi là Yên Ngựa khi nó lớn nhất trên dòng và nhỏ nhất trên cột
\end{problem}


\section{Kỹ thuật đặt cờ hiệu}
\begin{problem}
    Kiểm tra ma trận có tồn tại số dương hay không
\end{problem}
\begin{problem}
    Kiểm tra ma trận có tồn tại số hoàn thiện hay không
\end{problem}
\begin{problem}
    Kiểm tra ma trận có tồn tại số lẻ hay không
\end{problem}
\begin{problem}
    Kiểm tra ma trận có toàn dương hay không
\end{problem}
\begin{problem}
    Kiểm tra một hàng ma trận có tăng dần hay không
\end{problem}
\begin{problem}
    Kiểm tra một cột ma trận có giảm dần hay không. KiemTraCotGiamDan(a, n, m, k)
\end{problem}
\begin{problem}
    Kiểm tra các giá trị trong ma trận có giảm dần theo dòng và cột hay không
\end{problem}
\begin{problem}
    Liệt kê các dòng toàn âm trong ma trận các số thực
\end{problem}
\begin{problem}
    Liệt kê chỉ số các dòng chứa toàn giá trị chẵn trong ma trận các số nguyên
\end{problem}
\begin{problem}
    Liệt kê các dòng có chứa số nguyên tố trong ma trận các số nguyên
\end{problem}
\begin{problem}
    Liệt kê các dòng có chứa giá trị chẵn trong ma trận các số nguyên
\end{problem}
\begin{problem}
    Liệt kê các dòng có chứa giá trị âm trong ma trận các số thực
\end{problem}
\begin{problem}
    Liệt kê các cột trong ma trận các số nguyên có chứa số chính phương
\end{problem}
\begin{problem}
    Liệt kê các dòng trong ma trận các số thực thỏa mãn đồng thời các điều kiện sau:  dòng có chứa giá trị âm, giá trị 0 và giá trị dương
\end{problem}
\begin{problem}
    Liệt kê các dòng giảm dần trong ma trận
\end{problem}
\begin{problem}
    Liệt kê các cột tăng dần trong ma trận
\end{problem}
\begin{problem}
    Cho 2 ma trận A và B. Kiểm tra xem ma trận A có là ma trận con của ma trận B hay không

 1 2 3\\
 4 5 6\\

 6 7 3 1\\
 2 \textbf{1 2 3}\\
 5 \textbf{4 5 6}\\
 9 4 5 2\\
\end{problem}
\begin{problem}
    Cho 2 ma trận A và B. Đếm số lần xuất hiện của ma trận A trong ma trận B
\end{problem}


\section{Kỹ thuật đặt lính canh}
\begin{problem}
    Tìm số chẵn đầu tiên trong ma trận
\end{problem}
\begin{problem}
    Tìm max trong ma trận
\end{problem}
\begin{problem}
    Tìm giá trị lớn thứ 2 trong ma trận
\end{problem}
\begin{problem}
    Tìm số dương đầu tiên trong ma trận
\end{problem}
\begin{problem}
    Tìm giá trị âm lớn nhất trong ma trận
\end{problem}
\begin{problem}
    Liệt kê các dòng có chứa max
\end{problem}
\begin{problem}
    Tìm giá trị lớn nhất trên 1 dòng
\end{problem}
\begin{problem}
    Tìm giá trị nhỏ nhất trên 1 cột
\end{problem}
\begin{problem}
    Tìm số nguyên tố đầu tiên
\end{problem}
\begin{problem}
    Tìm số chẵn lớn nhất
\end{problem}
\begin{problem}
    Tìm số dương nhỏ nhất
\end{problem}
\begin{problem}
    Tìm số nguyên tố lớn nhất
\end{problem}
\begin{problem}
    Tìm 1 chữ số xuất hiện nhiều nhất
\end{problem}
\begin{problem}
     Đếm số lần xuất hiện của phần tử nhỏ nhất
\end{problem}
\begin{problem}
    Đếm số lượng chẵn nhỏ nhất
\end{problem}
\begin{problem}
    Tìm giá trị xuất hiện nhiều nhất
\end{problem}
\begin{problem}
    Tìm số chính phương lớn nhất
\end{problem}
\begin{problem}
    Tìm số hoàn thiện nhỏ nhất
\end{problem}
\begin{problem}
    Tìm các chữ số xuất hiện nhiều nhất trong ma trận
\end{problem}
\begin{problem}
    Liệt kê các dòng có tổng lớn nhất
\end{problem}
\begin{problem}
    Liệt kê các cột có tổng nhỏ nhất
\end{problem}
\begin{problem}
    Liệt kê các dòng có nhiều số chẵn nhất
\end{problem}
\begin{problem}
    Liệt kê các dòng có nhiều số nguyên tố nhất
\end{problem}
\begin{problem}
    Liệt kê các dòng có nhiều số hoàn thiện nhất
\end{problem}
\begin{problem}
    (*): Liệt kê các cột nhiều chữ số nhất trong ma trận các số nguyên
\end{problem}
\begin{problem}
    (*): Tìm ma trận con có tổng lớn nhất
\end{problem}


\section{Kỹ thuật xử lí ma trận}

\begin{problem}
    Hoán vị 2 dòng trên ma trận

1 2 3

4 5 6

7 8 9

Sau khi hoan vi dòng 0 và dòng 2 thì kết quả là

7 8 9 

4 5 6

1 2 3
\end{problem}
\begin{problem}
    Hoán vị 2 cột trên ma trận
\end{problem}
\begin{problem}
    Dịch xuống xoay vòng các hàng trong ma trận
\end{problem}
\begin{problem}
    Dịch lên xoay vòng các hàng trong ma trận
\end{problem}
\begin{problem}
    Dịch trái xoay vòng các cột trong ma trận
\end{problem}
\begin{problem}
    Dịch phải xoay vòng các cột trong ma trận
\end{problem}
\begin{problem}
    Dịch phải xoay vòng theo chiều kim đồng hồ các giá trị nằm trên biên ma trận
\end{problem}
\begin{problem}
    Dịch trái xoay vòng theo chiều kim đồng hồ các giá trị nằm trên biên ma trận
\end{problem}
\begin{problem}
    Xóa 1 dòng trong ma trận
\end{problem}
\begin{problem}
    Xóa 1 cột trong ma trận
\end{problem}
\begin{problem}
    Xoay ma trận 1 góc 90 độ
\end{problem}
\begin{problem}
    Xoay ma trận 1 góc 180 độ
\end{problem}
\begin{problem}
    Xoay ma trận 1 góc 270 độ
\end{problem}
\begin{problem}
    Chiếu gương ma trận theo trục dọc
\end{problem}
\begin{problem}
    Chiếu gương ma trận theo trục ngang
\end{problem}


\section{Kỹ thuật sắp xếp}
\begin{problem}
    Viết hàm sắp xếp các phần tử trên 1 dòng tăng dần từ trái sang phải
\end{problem}
\begin{problem}
    Viết hàm sắp xếp các phần tử trên 1 dòng giảm dần từ trái sang phải
\end{problem}
\begin{problem}
    Viết hàm sắp xếp các phần tử trên 1 cột tăng dần từ trên xuống dưới
\end{problem}
\begin{problem}
     Viết hàm sắp xếp các phần tử trên 1 cột giảm dần từ trên xuống dưới
\end{problem}
\begin{problem}
    Viết hàm xuất các giá trị chẵn trong ma trận các số nguyên theo thứ tự giảm dần
\end{problem}
\begin{problem}
    Viết hàm xuất các số nguyên tố trong ma trận các số nguyên theo thứ tự tăng dần
\end{problem}
\begin{problem}
    Viết hàm sắp xếp các phần tử trong ma trận theo yêu cầu sau:

    Dòng có chỉ số chẵn tăng dần

    Dòng có chỉ số lẻ giảm dần                            
\end{problem}
\begin{problem}
    Viết hàm sắp xếp các phần tử trong ma trận theo yêu cầu sau:

    Cột có chỉ số chẵn giảm dần từ trên xuống

    Cột có chỉ số lẻ tăng dần từ trên xuống
\end{problem}
\begin{problem}
    Sắp xếp ptử tăng dần theo hàng và cột: Dùng 2 phương pháp: sử dụng mảng phụ và ko dùng mảng phụ
\end{problem}
\begin{problem}
    Sắp xếp ptử dương tăng dần theo cột và dòng. Dùng 2 phương pháp: Sử dụng mảng phụ và ko sử dụng mảng phụ
\end{problem}
\begin{problem}
    Sắp xếp ptử chẵn giảm dần theo cột và dòng. Dùng 2 phương pháp: Sử dụng mảng phụ và ko sử dụng mảng phụ
\end{problem}
\begin{problem}
    Sắp xếp âm tăng dần, dương giảm dần, 0 giữ nguyên
\end{problem}
\begin{problem}
    Sắp xếp chẵn tăng, lẻ giảm
\end{problem}
\begin{problem}
    Sắp xếp các giá trị nằm trên biên ma trận tăng dần
\end{problem}
\begin{problem}
    Sắp xếp các giá trị dương nằm trên biên ma trận tăng dần
\end{problem}
\begin{problem}
    Sắp xếp các dòng dựa vào: tổng các ptử trong 1 dòng: sắp tăng dần.
\end{problem}
\begin{problem}
    Sắp xếp giá trị các ptử trong ma trận tăng dần theo dạng xoắn ốc (ma trận xoắn ốc)
\end{problem}
\begin{problem}
    Sắp xếp giá trị các ptử trong ma trận tăng dần theo dạng ziczac
\end{problem}
\begin{problem}
    Xuất các giá trị âm  giảm dần(ma trận không thay đổi sau khi xuất)
\end{problem}

\section{Kỹ thuật xây dựng ma trận}
\begin{problem}
    Cho ma trận A. Hãy tạo ma trận B, B[i][j] = abs(A[i][j])
\end{problem}
\begin{problem}
    Cho ma trận A. Hãy tạo ma trận B, B[i][j] = lớn nhất dòng i, cột j của A
\end{problem}
\begin{problem}
    Cho ma trận A. Hãy tạo ma trận B, B[i][j] = số lượng ptử dương xung quanh A[i][j]
\end{problem}
